
\begin{ex} % Câu 9
	Rút gọn biểu thức $B=\left(\frac{15-\sqrt{x}}{x-25}+\frac{2}{\sqrt{x}+5}\right):\frac{\sqrt{x}+3}{\sqrt{x}-5}$ với $x \ge 0, x \ne 25$.
	\loigiai{
		Điều kiện: $x \ge 0, x \ne 25$.
		Ta có $x-25 = (\sqrt{x}-5)(\sqrt{x}+5)$.
		\begin{align*}
			B &= \left(\frac{15-\sqrt{x}}{(\sqrt{x}-5)(\sqrt{x}+5)} + \frac{2(\sqrt{x}-5)}{(\sqrt{x}+5)(\sqrt{x}-5)}\right) : \frac{\sqrt{x}+3}{\sqrt{x}-5} \\
			&= \left(\frac{15-\sqrt{x} + 2\sqrt{x}-10}{(\sqrt{x}-5)(\sqrt{x}+5)}\right) \cdot \frac{\sqrt{x}-5}{\sqrt{x}+3} \\
			&= \frac{\sqrt{x}+5}{(\sqrt{x}-5)(\sqrt{x}+5)} \cdot \frac{\sqrt{x}-5}{\sqrt{x}+3} \\
			&= \frac{1}{\sqrt{x}-5} \cdot \frac{\sqrt{x}-5}{\sqrt{x}+3}
		\end{align*}
		Vì $x \ne 25$ nên $\sqrt{x}-5 \ne 0$. Ta có thể rút gọn:
		\begin{align*}
			B &= \frac{1}{\sqrt{x}+3}
		\end{align*}
		Vậy, với $x \ge 0, x \ne 25$ thì $B = \frac{1}{\sqrt{x}+3}$.
	}
\end{ex}

\begin{ex} % Câu 10
	Cho phương trình: $-3x^2+x+1=0$. Với $x_1, x_2$ là nghiệm của phương trình, không giải phương trình, hãy tính: $A=x_1^2+\frac{2}{x_1}+x_2^2+\frac{2}{x_2}$.
	\loigiai{
		Phương trình đã cho là $-3x^2+x+1=0$.
		Theo hệ thức Vi-ét, ta có:
		$ \begin{cases} x_1 + x_2 = -\frac{1}{-3} = \frac{1}{3} \\ x_1 x_2 = \frac{1}{-3} = -\frac{1}{3} \end{cases} $
		Ta biến đổi biểu thức A:
		\begin{align*}
			A &= x_1^2+\frac{2}{x_1}+x_2^2+\frac{2}{x_2} \\
			&= (x_1^2+x_2^2) + \left(\frac{2}{x_1}+\frac{2}{x_2}\right) \\
			&= ((x_1+x_2)^2 - 2x_1x_2) + \frac{2(x_1+x_2)}{x_1x_2}
		\end{align*}
		Thay giá trị từ hệ thức Vi-ét vào biểu thức A:
		\begin{align*}
			A &= \left(\left(\frac{1}{3}\right)^2 - 2\left(-\frac{1}{3}\right)\right) + \frac{2\left(\frac{1}{3}\right)}{-\frac{1}{3}} \\
			&= \left(\frac{1}{9} + \frac{2}{3}\right) + \frac{\frac{2}{3}}{-\frac{1}{3}} \\
			&= \left(\frac{1}{9} + \frac{6}{9}\right) - 2 \\
			&= \frac{7}{9} - 2 \\
			&= \frac{7}{9} - \frac{18}{9} \\
			&= -\frac{11}{9}
		\end{align*}
		Vậy, giá trị của biểu thức $A = -\frac{11}{9}$.
	}
\end{ex}

\begin{ex} % Câu 11
	Giải hệ phương trình sau: $ \begin{cases} -5x+2y=4 \\ 6x-3y=-7 \end{cases} $
	\loigiai{
		Ta có hệ phương trình:
		$ \begin{cases} -5x+2y=4 \quad &(1) \\ 6x-3y=-7 \quad &(2) \end{cases} $
		Nhân phương trình (1) với 3 và phương trình (2) với 2:
		$ \begin{cases} -15x+6y=12 \quad &(1') \\ 12x-6y=-14 \quad &(2') \end{cases} $
		Cộng vế theo vế của phương trình (1') và (2'):
		\begin{align*}
			(-15x+6y) + (12x-6y) &= 12 + (-14) \\
			-15x + 6y + 12x - 6y &= -2 \\
			-3x &= -2 \\
			x &= \frac{-2}{-3} \\
			x &= \frac{2}{3}
		\end{align*}
		Thay $x=\frac{2}{3}$ vào phương trình (1):
		\begin{align*}
			-5\left(\frac{2}{3}\right) + 2y &= 4 \\
			-\frac{10}{3} + 2y &= 4 \\
			2y &= 4 + \frac{10}{3} \\
			2y &= \frac{12}{3} + \frac{10}{3} \\
			2y &= \frac{22}{3} \\
			y &= \frac{22}{3 \cdot 2} \\
			y &= \frac{11}{3}
		\end{align*}
		Vậy, nghiệm của hệ phương trình là $(x;y) = \left(\frac{2}{3};\frac{11}{3}\right)$.
	}
\end{ex}

\begin{ex} % Câu 12a - Xác suất gieo xúc xắc
	Bạn An chơi trò chơi gieo đồng thời hai con xúc xắc cân đối đồng chất I và II để xem qua các lần gieo xúc xắc mình có thể gieo nhiều nhất được bao nhiêu điểm. Tính xác suất của các biến cố A: “Có đúng 1 con xúc xắc xuất hiện mặt 6 chấm”.
	\loigiai{
		Khi gieo đồng thời hai con xúc xắc, không gian mẫu $\Omega$ có $6 \times 6 = 36$ phần tử.
		Biến cố A: “Có đúng 1 con xúc xắc xuất hiện mặt 6 chấm”.
		Các trường hợp thuận lợi cho biến cố A:
		\begin{itemize}
			\item Xúc xắc I ra mặt 6 chấm, xúc xắc II không ra mặt 6 chấm (ra 1, 2, 3, 4, hoặc 5): có $1 \times 5 = 5$ trường hợp.
			(Các cặp: (6,1), (6,2), (6,3), (6,4), (6,5))
			\item Xúc xắc I không ra mặt 6 chấm, xúc xắc II ra mặt 6 chấm: có $5 \times 1 = 5$ trường hợp.
			(Các cặp: (1,6), (2,6), (3,6), (4,6), (5,6))
		\end{itemize}
		Tổng số trường hợp thuận lợi cho biến cố A là $n(A) = 5 + 5 = 10$.
		Xác suất của biến cố A là:
		$P(A) = \frac{n(A)}{n(\Omega)} = \frac{10}{36} = \frac{5}{18}$.
	}
\end{ex}

\begin{ex} % Bài toán hội trường (1 ẩn)
	Một hội trường có 300 ghế ngồi được xếp đều thành từng dãy (số ghế trên mỗi dãy bằng nhau). Nếu kê thêm 2 dãy và mỗi dãy xếp thêm 1 ghế thì trong hội trường có tất cả 351 ghế. Hỏi lúc đầu hội trường có bao nhiêu dãy ghế.
	\loigiai{
		Gọi $x$ là số dãy ghế lúc đầu của hội trường.
		Điều kiện: $x$ là số nguyên dương.
		
		Lúc đầu, hội trường có 300 ghế, vậy số ghế trên mỗi dãy lúc đầu là: $\frac{300}{x}$ (ghế).
		
		Sau khi thay đổi:
		Số dãy ghế mới là: $x+2$ (dãy).
		Số ghế trên mỗi dãy mới là: $\frac{300}{x} + 1$ (ghế).
		Tổng số ghế mới trong hội trường là 351 ghế, nên ta có phương trình:
		\begin{align*}
			(x+2)\left(\frac{300}{x}+1\right) &= 351
		\end{align*}
		Để giải phương trình này, ta tiến hành biến đổi:
		\begin{align*}
			(x+2)\left(\frac{300+x}{x}\right) &= 351 \\
			(x+2)(300+x) &= 351x \quad (\text{Nhân cả hai vế với } x, \text{ vì } x>0) \\
			300x + x^2 + 600 + 2x &= 351x \\
			x^2 + 302x + 600 &= 351x \\
			x^2 + 302x - 351x + 600 &= 0 \\
			x^2 - 49x + 600 &= 0
		\end{align*}
		Giải phương trình bậc hai $x^2 - 49x + 600 = 0$:
		$\Delta = b^2 - 4ac = (-49)^2 - 4(1)(600)$
		\begin{align*}
			\Delta &= 2401 - 2400 \\
			\Delta &= 1
		\end{align*}
		$\sqrt{\Delta} = \sqrt{1} = 1$.
		Phương trình có hai nghiệm phân biệt:
		$x_1 = \frac{-(-49) + 1}{2 \cdot 1} = \frac{49 + 1}{2} = \frac{50}{2} = 25$
		$x_2 = \frac{-(-49) - 1}{2 \cdot 1} = \frac{49 - 1}{2} = \frac{48}{2} = 24$
		
		Cả hai giá trị $x=25$ và $x=24$ đều là số nguyên dương, nên ta cần kiểm tra xem giá trị nào làm cho số ghế trên mỗi dãy là số nguyên.
		
		Trường hợp 1: Nếu $x = 25$ (số dãy ghế lúc đầu là 25)
		Số ghế trên mỗi dãy lúc đầu là: $\frac{300}{25} = 12$ (ghế). Đây là số nguyên, hợp lệ.
		Kiểm tra lại:
		Số dãy mới: $25+2 = 27$ dãy.
		Số ghế mỗi dãy mới: $12+1 = 13$ ghế.
		Tổng số ghế mới: $27 \times 13 = 351$ ghế (Đúng với đề bài).
		
		Trường hợp 2: Nếu $x = 24$ (số dãy ghế lúc đầu là 24)
		Số ghế trên mỗi dãy lúc đầu là: $\frac{300}{24} = 12,5$ (ghế). Đây không phải là số nguyên, nên trường hợp này không hợp lệ vì số ghế trên mỗi dãy phải là số nguyên.
		
		Vậy, lúc đầu hội trường có 25 dãy ghế.
	}
\end{ex}
\begin{ex} % Câu 13
	Cho đường tròn $(O;R)$ và dây $MN$ cố định ($MN < 2R$). Kẻ đường kính $AB$ vuông góc với dây $MN$ tại $E$ (điểm $A$ thuộc cung nhỏ $MN$). Lấy điểm $C$ thuộc dây $MN$ ($C$ khác $M, N, E$). Đường thẳng $BC$ cắt $(O;R)$ tại điểm $K$ ($K$ khác $B$).
	\begin{enumerate}[a)]
		\item Chứng minh $AKCE$ là tứ giác nội tiếp.
		\item Chứng minh $BM^2 = BK \cdot BC$. Gọi $I$ là giao điểm của hai đường thẳng $AK$ và $MN$; $D$ là giao điểm của hai đường thẳng $AC$ và $BI$. Chứng minh điểm $C$ cách đều ba cạnh của $\triangle DEK$.
	\end{enumerate}
	\loigiai{}
\end{ex}


\setcounterref{ex}{0}
